% !PRINTED 223
\weberpage{223}

%<<CONT
and thus $\mathrm{sn}\,\Omega$ any root of the proper division equation, then the sum, to be
extended over all $\alpha$,
\weq{13}{\Sigma(\pm 1)^\varepsilon(\pm 1)^{\varepsilon'}\ldots\, \mathrm{sn}\,\alpha\Omega = \psi(\Omega),}
in which the $k$ signs of $\pm 1$ may be chosen arbitrarily, has the property that for every
$\alpha$ determined according to (12)
\weq{14}{\psi(\alpha\Omega) = (\pm 1)^\varepsilon(\pm 1)^{\varepsilon'}\ldots \psi(\Omega)}
holds; and the square of $\psi(\Omega)$ remains unchanged under the substitutions of the group
$\mathfrak{B}_0$, and is therefore rationally expressible by $n$th roots of unity and the roots
of a transformation equation. $\psi(\Omega)$ is therefore the square root $\sqrt{A}$ of such an
expression $A$. But the number of the expressions (13) amounts to $2^k$, and by addition of all
the equations so formed we obtain
\weq{15}{2^k\,\mathrm{sn}\,\Omega = \Sigma\sqrt{A}.}

It is still to be remarked that of the quantities $\psi$, and hence also of the $A$, half
vanish. For we have
\[
-1 \equiv cc' \ldots \ (\mathrm{mod}\ n),
\]
and if we therefore set $\alpha \equiv -1$ in (14), it follows that:
\[
\psi(-\Omega) = (\pm 1)(\pm 1)\ldots\psi(\Omega);
\]
but on the other hand
\[
\psi(-\Omega) = -\psi(\Omega),
\]
and consequently $\psi(\Omega)$ always vanishes when among the $k$ quantities $\pm 1$ occurring
in (14) an even number of minus signs occurs.

The roots $x_{\mu,\mu'}$ can therefore be expressed linearly by $2^{k-1}$ square roots.

\gesperrt{If $n$ is a prime or a power of a prime, then the squares of the roots of the division equation
are rationally expressible by the roots of the transformation equation}\footnote{Cf. on this
theorem \gesperrt{Kronecker}, Monatsbericht der Berliner Akademie, 19.\ Juli 1875. There \gesperrt{Kronecker} draws
attention to an oversight which occurs in a memoir of \gesperrt{Jacobi} bearing on this subject (Crelle,
Bd.\ 47 and Bd.\ 50).}.

If $n$ is a composite number, the classification of the roots $x_{\mu,\mu'}$ into rows can be
carried still further. If $p$ is a prime divisor of
%CONT>>
